\section{Entry, shell summation, and the limiting distance measure}
\label{asm:section}
\label{sec:assembly}

The graph estimate and the two-depth potential now give decay on each frequency
shell.  The measures used on different shells have different classes and
masks.  We first explain how to restore their mass, and then give a scalar
reconstruction argument that accommodates the decreasing pair filters.
Throughout this Section all distances are measured in the original
space $\R^d$.

\subsection{The entry depth}

Let $S=d/2$.  Fix the starting probabilities $\mu_1,\mu_2$ and the number
$\chi>0$ supplied by Proposition~\ref{prep:starting-pairs}.
Make one common similarity placing $E$ in $(1/4,3/4)^d$, and retain
the notation for the transformed sets and measures.  This multiplies
every distance by one fixed positive constant and preserves the
desired positive-measure conclusion.  Choose one small common
translation of the dyadic grid whose boundaries have zero mass for
both starting laws.  Its depth-zero cube $Q_0$ still contains $E$ in
its interior.  Use this same grid for both class decompositions.
The two actual supports remain separated by a fixed positive
distance, and their cross directions remain in the prepared patch.
At scale
$R=2^N$, apply Lemma~\ref{reg:regularization} with parameter $\sigma$ to
both probabilities.  Write $\mathcal C_{j,N}$ for the retained classes on
side $j$.  For $C\in\mathcal C_{j,N}$, let
\[
 a_{j,C}=\mu_j(C),\qquad
 \rho_{j,C}=a_{j,C}^{-1}\mu_j|_C.
\]
The classes on each side are disjoint, their number is $R^{o(1)}$, and
their discarded original mass is at most $R^{-\sigma+o(1)}$.  All the
$o(1)$ terms below are uniform over these classes, with the fixed
parameters held constant.
Every normalized class is supported in the same root cube $Q_0$ and
assigns it mass one.  Since every profile has $m_j(0)=0$, its root
normalization is exactly $M_j(0)=R^{-m_j(0)}=1$.

For a pair of classes, suppress their labels and write
\[
 A_{j,n}=m_j(n/N)-Sn/N,\qquad
 k_j(g,p)=A_{j,p}-\min_{g\le n\le p}A_{j,n}.
\]
Recall from Lemma~\ref{test:profile-bounds} that
\begin{equation}
 \label{asm:profile-input}
 A_{j,0}=0,\qquad
 |A_{j,n}-A_{j,n'}|\le d|n-n'|/N,\qquad
 A_{j,n}\ge\chi n/N-C_d\sigma.
\end{equation}
Here and in what follows, constants denoted by $C_d$ may be enlarged a
finite number of times and depend only on $d$.

We seek a common depth $c$ for which one entry cost $k_j(0,c)$ is small
and both excess profiles remain above a positive level on $[c,N]$.
The graph estimate contributes the smaller cost.  After
Cauchy--Schwarz separates the profiles, the two single-profile
estimates must supply a larger potential saving.

\begin{lemma}[Entry depth]\label{asm:entry}
Fix $0<\beta<\min(\chi/10,1/10)$.  If $\sigma$ is sufficiently small
in terms of $d,\chi,\beta$, then, for all sufficiently large $N$, the
integer
\begin{equation}
 \label{asm:entry-definition}
 c=\max\{n\in\{0,\ldots,N\}:
                  \min(A_{1,n},A_{2,n})\le\beta\}
\end{equation}
is defined and satisfies
\begin{align}
 \label{asm:entry-range}
 &\left\lfloor\frac{\beta N}{2d}\right\rfloor\le c<N/4,\\
 \label{asm:entry-lower}
 &A_{j,n}\ge\beta-d/N
       \quad(j=1,2,\ c\le n\le N),\\
 \label{asm:entry-cost}
 &\min_{j=1,2} k_j(0,c)\le\beta+C_d\sigma.
\end{align}
If $\kappa$ is sufficiently small and $q_0=\lceil\kappa N\rceil$, the
interval $[c,N]$ is nonterminal.  Its potentials from
Definition~\ref{ledger:potential} obey
\begin{equation}
 \label{asm:entry-potential}
 \mathcal V_j([c,N])
 :=\min_{\substack{c\le u<v\le N\\v-u\ge10q_0}}
                         (A_{j,u}+A_{j,v})
 \ge2\beta-2d/N.
\end{equation}
\end{lemma}

\begin{proof}
The index $0$ qualifies in Equation~\eqref{asm:entry-definition}.
For $n\le\beta N/(2d)$, the Lipschitz bound in
Equation~\eqref{asm:profile-input} gives $A_{j,n}\le\beta/2$, proving
the lower bound in Equation~\eqref{asm:entry-range}.  For $n\ge N/4$,
the lower profile bound gives
$A_{j,n}\ge\chi/4-C_d\sigma>\beta$ when $\sigma$ is sufficiently
small.  This proves its upper bound.

By maximality of $c$, both $A_{j,c+1}$ exceed $\beta$, so
$A_{j,c}>\beta-d/N$.  Both profiles exceed $\beta$ at every later
index.  For at least one side $j_*$, we also have
$A_{j_*,c}\le\beta$.  Since
$\min_{0\le n\le c}A_{j_*,n}\ge-C_d\sigma$, this gives
Equation~\eqref{asm:entry-cost}.  The width of $[c,N]$ exceeds $3N/4$;
thus it exceeds $100q_0$ for sufficiently small fixed $\kappa$ and
large $N$.  Every pair entering its potential consists of two indices
covered by Equation~\eqref{asm:entry-lower}, which proves
Equation~\eqref{asm:entry-potential}.
\end{proof}

\subsection{Positive approximants and restoration of mass}

Write
\[
 D(x,y)=|x-y|,\qquad e(x,y)=\frac{x-y}{|x-y|},\qquad
 W(x,y)=D(x,y)^{-(d-1)/2}.
\]
The fixed separated supports give constants $0<d_0<d_1<\infty$ with
\begin{equation}
 \label{asm:distance-range}
 d_0\le D(x,y)\le d_1
 \quad\hbox{on }\supp\mu_1\times\supp\mu_2.
\end{equation}
In particular $W$ is bounded above and below there by positive constants.

For each pair of classes $(C_1,C_2)$, use its own entry index $c$ from
Equation~\eqref{asm:entry-definition}.  At endpoint $j$ take the extra
parallel descriptor $[0,c]$, together with every descriptor in
$\mathcal M_j(c,N)$.  Let
\begin{equation}
 \label{asm:endpoint-mask}
 b_{j;C_1,C_2}(z,w)
 =w^{(j)}_{0,[0,c]}(z,w)
                  \prod_{\mathfrak t\in\mathcal M_j(c,N)}
                         w^{(j)}_{0,\mathfrak t}(z,w).
\end{equation}
The masks are those of Lemma~\ref{test:smooth-masks}; thus every factor
lies in $[0,1]$ and equals one when its level-zero condition holds.
Antipodal symmetry allows both endpoints to be evaluated at $e(x,y)$.
Define the Borel function $\Theta_N$ on the fixed separated product
$E_1\times E_2$ by
\[
 \Theta_N(x,y)=
 b_{1;C_1,C_2}(x,e(x,y))
 b_{2;C_1,C_2}(y,e(x,y))
 \quad\hbox{if }(x,y)\in C_1\times C_2,
\]
and set it equal to zero whenever either endpoint belongs to no retained
class.
In particular $0\le\Theta_N\le1$.

Let $\Gamma_N$ be the decreasing pair filters of
Proposition~\ref{prep:starting-pairs}.  With $D_\#$ denoting pushforward
under the scalar map $D$, define
\begin{equation}
 \label{asm:measures}
 \alpha_N=D_\#\bigl(W\one_{\Gamma_N}
                          (\mu_1\otimes\mu_2)\bigr),\qquad
 \tau_N=D_\#\bigl(W\one_{\Gamma_N}\Theta_N
                          (\mu_1\otimes\mu_2)\bigr).
\end{equation}
These measures use the original class masses.  Only $\tau_N$ depends on
the shell's class decomposition and angular masks.

\begin{lemma}[Restoration of mass]\label{asm:retained-measures}
There is $\eta>0$, depending on the fixed parameters and starting
measures, such that
\begin{equation}
 \label{asm:tv-approximation}
 0\le\tau_N\le\alpha_N,\qquad
 \|\alpha_N-\tau_N\|_{\TV}\lesssim2^{-\eta N}.
\end{equation}
Moreover $\alpha_N$ decreases to the nonzero finite positive measure
\begin{equation}
 \label{asm:limit-measure}
 \alpha_\infty=D_\#\bigl(W\one_{\Gamma_\infty}
                               (\mu_1\otimes\mu_2)\bigr),\qquad
 \Gamma_\infty=\bigcap_N\Gamma_N.
\end{equation}
All these measures are carried by the compact set of original distances
$\Delta(E)$.
\end{lemma}

\begin{proof}
Apply the cap preparation in Proposition~\ref{prep:starting-pairs} with
loss $R^\sigma$.  It gives $\mathcal G_N\subset\Gamma_N$ such that
\[
 (\mu_1\otimes\mu_2)(\Gamma_N\setminus\mathcal G_N)
       \lesssim R^{-c_*}
 \quad\hbox{for some }c_*>0.
\]
For $\mu_1$-almost every first endpoint $x$, the angular subprobability
\[
 \xi_{N,x}=e(x,\cdot)_*
             \bigl(\one_{\mathcal G_N}(x,\cdot)\mu_2\bigr)
\]
has the required $S$-cap bound with loss $R^\sigma$ down to scale
$R^{-1}$.  The change from $\pi_x(y)$ to $e(x,y)=-\pi_x(y)$ preserves
cap bounds, since the antipodal map sends caps to caps.  The analogous
assertion holds for $\mu_2$-almost every second endpoint with the
endpoints reversed.
These are bounds for restrictions of the original opposite probability.

Fix $C_1,C_2$ and one descriptor $\mathfrak t$ used at the first endpoint.
For $\mu_1$-almost every $x\in C_1$, the anchor cell determined by
$\mathfrak t$ and $x$ is occupied and $\xi_{N,x}$ has the cap bound above.  Let
$F_{C_1,C_2,\mathfrak t,x}$ be the directions failing this level-zero test.
For such $x$, Lemma~\ref{test:exceptional-directions} applies directly
to $\xi_{N,x}$ and gives
\begin{equation}
 \label{asm:original-fiber-estimate}
 \mu_2\{y\in C_2:(x,y)\in\mathcal G_N,
       \ e(x,y)\in F_{C_1,C_2,\mathfrak t,x}\}\le R^{-\sigma}.
\end{equation}
Restricting the set of $y$ to $C_2$ only decreases the left side; no
normalization by $\mu_2(C_2)$ occurs.  Integrating in $x$ over $C_1$
costs at most $a_{1,C_1}R^{-\sigma}$.  Sum first over $C_1$, whose
original masses have sum at most one, and then over $C_2$ and the
descriptors.  The latter two counts are $R^{o(1)}$, by
Lemma~\ref{reg:regularization} and
Lemma~\ref{test:descriptor-properties}.  Thus the first-endpoint
failure mass is $R^{-\sigma+o(1)}$.  The same argument in the other
orientation gives this bound for second-endpoint failures.  There is
no sum over anchor cells: a point and descriptor already specify one.

The product in Equation~\eqref{asm:endpoint-mask} equals one if all
its level-zero tests hold.  Hence $1-\Theta_N$ on retained class pairs is
bounded by the indicator of the union of these failures.  Including
the original mass outside retained class pairs and
$\Gamma_N\setminus\mathcal G_N$ gives
\[
 \int_{\Gamma_N}(1-\Theta_N)\,d(\mu_1\otimes\mu_2)
 \lesssim R^{-c_*}+R^{-\sigma+o(1)}.
\]
The bounded weight $W$ proves Equation~\eqref{asm:tv-approximation},
for example with any fixed
$0<\eta<\min(c_*,\sigma)/2$ and sufficiently large $N$.
The set $\mathcal G_N$ was used only in this estimate: it is not a
multiplier in Equation~\eqref{asm:measures}.

The filters $\Gamma_N$ decrease and have positive limiting product
mass.  Continuity of a finite measure along decreasing sets proves
Equation~\eqref{asm:limit-measure} and monotonicity of $\alpha_N$.
The positive lower bound for $W$ shows
$\alpha_\infty(\R)>0$.  Finally, the two original supports are compact
subsets of $E$, and $D$ is continuous.  Every pushforward above is
therefore carried by their compact distance image, which is contained
in $\Delta(E)$.
\end{proof}

\subsection{The shell estimate}

\begin{proposition}[Decay of the approximating distance measures]
\label{asm:shell-decay}
The fixed parameters can be chosen, in the order
$\beta,\kappa,\sigma,\omega$, so that some $\gamma>0$ satisfies
\begin{equation}
 \label{asm:shell-bound}
 \int_{2^{N-1}\le |r|\le2^{N+1}}
                   |\widehat{\tau_N}(r)|^2\,dr
       \lesssim2^{-\gamma N}
\end{equation}
for all sufficiently large $N$.
\end{proposition}

\begin{proof}
For the moment fix one normalized class pair and one product
$X_\ell\times Y_\ell$ in the spatial partition of $\Gamma_N$.
The corresponding component measure $\nu$ is the pushforward of
\[
 W(x,y)\one_{X_\ell}(x)\one_{Y_\ell}(y)
 b_{1;C_1,C_2}(x,e(x,y))b_{2;C_1,C_2}(y,e(x,y))
             \,d\rho_{1,C_1}(x)\,d\rho_{2,C_2}(y).
\]
The product indicators are permitted bounded measurable spatial
multipliers in the oscillatory estimates.  At depth zero take
$X=Y=Q_0$, each with its own normalized class law.  These are occupied
root cells, with $M_1(0)=M_2(0)=1$.  Their supported laws lie in the
fixed separated sets $E_1,E_2$, so the outer-support geometry of
Section~\ref{osc:section} applies.  The cell-center separation in
Equation~\eqref{osc:separation} is the internal-cell convention; the
outer root uses the separation of the supported laws.  In particular,
all root spatial diameters in the graph argument are $O_d(1)$, while
the actual cross distances have a fixed positive lower bound.
By
Lemma~\ref{test:descriptor-properties}, the largest nominal length
of an endpoint descriptor is at most
\[
 L\le\max\bigl(c,(N-c)/2\bigr)\le N/2.
\]
Thus the root shell conversion, Proposition~\ref{osc:shell-conversion},
applies with $a=0$, $v=N$.  Its weight is precisely $W$; there is no
internal normalization of the distance weight at this root.  For any
prescribed fixed $T>0$, it bounds the squared Fourier norm on the positive
shell by $R^{o(1)}\sup\mathsf H_F(0,N)+O(R^{-T})$.  The supremum retains
the root cells, class normalizations, lists, and smoothing index zero;
its admissible symbols have the fixed derivative budget and common
numerical angular bounds supplied by that Proposition for $T$.

Apply Corollary~\ref{graph:separated-form} at $p=c$ to each energy in this
supremum.  The required parallel
descriptor is the extra $[0,c]$ on each side, with $g=0$.  All scale
conditions hold: $c>0$ for large $N$, $c<N/4$,
$L\le N/2<N-c$, and $N-c\ge10\sigma N$.  Since $A_{j,0}=0$, the
graph factor is
\begin{equation}
 \label{asm:graph-factor}
 R^{-A_{1,c}-A_{2,c}+\min_j k_j(0,c)+O_d(\sigma)+o(1)}.
\end{equation}
The graph estimate produces a general two-profile energy
$\mathsf H(c,N)$ on a pair of occupied $c$-cells, still carrying the
same endpoint symbols.

On these cells the extra $[0,c]$ factor is spatially constant within
each endpoint cell.  It, and any of its normalized differentiated
versions produced by shell conversion, can be removed at bounded
cost from the positive independent angular energies.  The remaining
descriptor lists are exactly those with base $[c,N]$ and with reached
anchors omitted.  This removal is made in the general positive energy
$\mathsf H$, where the two angular factors are nonnegative.

To separate the two profiles, write their nonnegative angular
energies as $E_1(r),E_2(r)$ and use the common measure
$d\lambda_N(r)=\Psi(r/2^N)\,dr/2^N$.  Cauchy--Schwarz gives the exact
normalization
\begin{equation}
 \label{asm:two-profile-cs}
 \begin{split}
 2^{d(N-c)}\int E_1(r)E_2(r)\,d\lambda_N(r)
 &\le
 \left(2^{d(N-c)}\int E_1(r)^2\,d\lambda_N(r)\right)^{1/2}\\
 &\quad\cdot
 \left(2^{d(N-c)}\int E_2(r)^2\,d\lambda_N(r)\right)^{1/2}.
 \end{split}
\end{equation}
Each parenthesis is the independent angular energy formed by
duplicating one endpoint law, cell, and symbol.  With the unreached
descriptors of $\mathcal M_j(c,N)$, it is the single-profile recursive
state $\mathsf H_0^{[c,N]}(c,N)$ for profile $j$; denote this value by
$\mathsf H_j(c,N)$.  Its symbols have the fixed derivative budget and
common numerical angular bounds supplied by shell conversion;
restriction to the $c$-cells and duplication preserve these bounds.
These states impose no separation condition, so
Theorem~\ref{ledger:bound} applies to each, including the case of two
identical cells.  It yields
\[
 \mathsf H_j(c,N)
 \le R^{2A_{j,c}-\mathcal V_j([c,N])
                         +O_d(\kappa+\sigma/\kappa)+o(1)}.
\]
The profile terms in Equation~\eqref{asm:graph-factor} and the
geometric mean of these bounds combine to
\begin{equation}
 \label{asm:entry-cancellation}
 \begin{split}
 &-A_{1,c}-A_{2,c}+\min_j k_j(0,c)
       +\frac12\sum_{j=1}^2
                    \bigl(2A_{j,c}-\mathcal V_j([c,N])\bigr)\\
 &\hspace{20mm}
 =\min_j k_j(0,c)-\frac12
                     \bigl(\mathcal V_1([c,N])+\mathcal V_2([c,N])\bigr)
 \le-\beta+C_d\sigma+2d/N,
 \end{split}
\end{equation}
where the last inequality uses the cost and potential bounds in
Lemma~\ref{asm:entry}.  Including the stated losses proves uniformly
for each component that
\begin{equation}
 \label{asm:component-shell}
 \int_{2^{N-1}\le r\le2^{N+1}}|\widehat\nu(r)|^2\,dr
 \lesssim
 R^{-\beta+C_d(\kappa+\sigma/\kappa)+o(1)}+R^{-T}.
\end{equation}
Here $T$ can be any fixed sufficiently large number, after choosing
the expansion and integration-by-parts orders.  The term $O_d(\sigma)$
has been included in $O_d(\sigma/\kappa)$, with $\kappa<1$.

We make the parameter and error choices explicit.  First choose
$\beta$ as in Lemma~\ref{asm:entry}.  Next choose $\kappa$ so small
that the entry interval is nonterminal and
$C_d\kappa<\beta/16$.  Then choose $\sigma$ so small that all the
profile and transition hypotheses hold and
$C_d\sigma/\kappa<\beta/16$.  Finally choose $\omega>0$ so small
that
\begin{equation}
 \label{asm:omega-choice}
 4d\omega<\beta/8.
\end{equation}
The preparation of $\Gamma_N$ is used with this fixed $\omega$ and
cap loss $R^\sigma$.  Its deletion exponent $c_*$ may depend on both
parameters; only $c_*>0$ was needed in
Lemma~\ref{asm:retained-measures}.

For completeness, the error accuracy is fixed after these parameters
and before $N$ tends to infinity.  Theorem~\ref{ledger:bound} bounds
the number of steps on a branch by a constant depending on
$\kappa,\sigma$.  Choose a fixed local negative-power accuracy large
enough to dominate all ancestor factors and all errors on such a
branch, leaving the root remainder $R^{-T}$ in
Equation~\eqref{asm:component-shell}, with, say, $T>10+4d\omega$.
The inverse expansion orders are then fixed and determine a finite
symbol derivative budget.  The direct expansion and
integration-by-parts orders are chosen for that budget and accuracy,
as in Theorem~\ref{ledger:bound}.  The graph step and
Equation~\eqref{asm:two-profile-cs} introduce no persistent new
derivative orders; the spatial indicators have zero angular
derivatives.  These choices affect constants and the threshold for
$N$, and do not change the dimensional coefficients in the exponent
losses above.  We now take $N$ sufficiently large for all these fixed
choices and for the uniform $o(1)$ losses.

There are at most $L_N\lesssim R^{2d\omega}$ spatial products in
$\Gamma_N$.  Let $\nu_{C_1,C_2,\ell}$ denote the preceding component.
The exact recombination is
\[
 \tau_N=\sum_{C_1,C_2}a_{1,C_1}a_{2,C_2}
                    \sum_{\ell=1}^{L_N}\nu_{C_1,C_2,\ell}.
\]
Triangle inequality in the Fourier $L^2$ norm costs no class-count
factor, because
\[
 \sum_{C_1,C_2}a_{1,C_1}a_{2,C_2}
 =\left(\sum_{C_1}a_{1,C_1}\right)
  \left(\sum_{C_2}a_{2,C_2}\right)\le1.
\]
It costs at most $L_N$ for the spatial products.  Squaring therefore
costs $R^{4d\omega}$, giving
\begin{equation}
 \label{asm:assembled-shell}
 \int_{2^{N-1}\le r\le2^{N+1}}
                  |\widehat{\tau_N}(r)|^2\,dr
 \lesssim R^{-\beta+C_d(\kappa+\sigma/\kappa)+4d\omega+o(1)}
                       +R^{-T+4d\omega}.
\end{equation}
Our successive choices leave a strictly negative exponent; for
example, increasing the starting $N$ if necessary, the right side is
$O(R^{-\beta/2})$.  Each $\tau_N$ is a real positive measure, so
$\widehat{\tau_N}(-r)=\overline{\widehat{\tau_N}(r)}$.  The same
bound holds on the negative shell, proving
Equation~\eqref{asm:shell-bound} with some $\gamma>0$.
\end{proof}

\subsection{A monotone shell reconstruction lemma}

The following scalar fact isolates the final limiting argument.  Its
hypotheses require no rate at which the decreasing measures approach
their limit.  The proof combines a smooth dyadic Fourier decomposition
with the summable positive increments of a decreasing family of finite
measures. We give this reconstruction explicitly.

\begin{lemma}[Monotone measure shell criterion]\label{asm:monotone-shell}
Let $(\alpha_N)_{N\ge K}$ be finite positive Borel measures on $\R$
with $\alpha_{N+1}\le\alpha_N$, and let
$\alpha_\infty=\lim_N\alpha_N$.  Let $\tau_N$ be finite complex
Borel measures such that
\begin{equation}
 \label{asm:criterion-hypotheses}
 \sum_{N\ge K}\|\alpha_N-\tau_N\|_{\TV}<\infty,
 \qquad
 \sum_{N\ge K}
 \left(\int_{2^{N-1}\le |r|\le2^{N+1}}
                   |\widehat{\tau_N}(r)|^2\,dr\right)^{1/2}<\infty.
\end{equation}
Then $\alpha_\infty$ is absolutely continuous with respect to
Lebesgue measure.  In particular, the conclusion holds if the two
quantities inside these sums are respectively
$O(2^{-\eta N})$ and $O(2^{-\gamma N/2})$ for arbitrary
$\eta,\gamma>0$.
\end{lemma}

\begin{proof}
Choose $\vartheta\in C_c^\infty(\R)$ equal to one on $[-1,1]$ and
supported in $[-2,2]$.  Let $Q_N$ be convolution with the inverse
Fourier transform of $\vartheta(2^{-N}r)$, and put
$P_N=Q_N-Q_{N-1}$.  The multiplier of $P_N$ is supported in
$\{2^{N-1}\le|r|\le2^{N+1}\}$.  Dilation of the fixed Schwartz
kernel gives a constant $C_\vartheta$ such that, for every finite
measure $\nu$,
\begin{equation}
 \label{asm:low-cutoff-bound}
 \|Q_N\nu\|_{L^1}\le C_\vartheta\|\nu\|_{\TV},\qquad
 \left\|\sum_{N=a}^bP_N\nu\right\|_{L^1}
 \le2C_\vartheta\|\nu\|_{\TV}\quad(a\le b).
\end{equation}
The second bound follows from telescoping to $Q_b-Q_{a-1}$; in
particular its constant is independent of the number of bands.

With the convention $\widehat\nu(r)=\int e^{-irt}\,d\nu(t)$,
Plancherel \cite[equation~(12.4)]{Mattila1995} and the uniformly
bounded band multipliers give
\[
 \|P_N\tau_N\|_{L^2}
 \lesssim
 \left(\int_{2^{N-1}\le |r|\le2^{N+1}}
                   |\widehat{\tau_N}(r)|^2\,dr\right)^{1/2}.
\]
Thus
\begin{equation}
 \label{asm:two-convergent-series}
 G:=\sum_{N\ge K}P_N\tau_N\quad\hbox{converges in }L^2,
 \qquad
 H:=\sum_{N\ge K}P_N(\alpha_N-\tau_N)
                         \quad\hbox{converges in }L^1.
\end{equation}
Both series converge absolutely in the indicated Banach space.

It remains to account for the variation of $\alpha_N$.  Set
$\zeta_n=\alpha_n-\alpha_{n+1}\ge0$.  Monotonicity and finiteness
give
\begin{equation}
 \label{asm:measure-telescoping}
 \sum_{n\ge K}\|\zeta_n\|_{\TV}
       =\alpha_K(\R)-\alpha_\infty(\R)<\infty,
 \qquad
 \alpha_N-\alpha_\infty=\sum_{n\ge N}\zeta_n
                         \quad\hbox{in total variation}.
\end{equation}
For a finite $M\ge K$, define
$Z_M=\sum_{N=K}^MP_N(\alpha_N-\alpha_\infty)$.  Substituting
Equation~\eqref{asm:measure-telescoping} into this finite sum and
using the bounded operators from measures to $L^1$ gives the exact
identity
\begin{equation}
 \label{asm:finite-band-identity}
 Z_M=\sum_{n\ge K}
             (Q_{\min(n,M)}-Q_{K-1})\zeta_n.
\end{equation}
For each fixed $M$ the series on the right converges in $L^1$.
Moreover, Equation~\eqref{asm:low-cutoff-bound} shows that
\[
 Z:=\sum_{n\ge K}(Q_n-Q_{K-1})\zeta_n
\]
converges absolutely in $L^1$.  Subtracting the two series cancels
all $n\le M$, and gives
\begin{equation}
 \label{asm:variation-tail}
 \begin{split}
 \|Z_M-Z\|_{L^1}
 &\le2C_\vartheta\sum_{n>M}\|\zeta_n\|_{\TV}\\
 &=2C_\vartheta
             \|\alpha_{M+1}-\alpha_\infty\|_{\TV}
 \longrightarrow0.
 \end{split}
\end{equation}
This is where monotonicity is used.  The original double series has
not been rearranged without a convergence justification.

Finally, finite telescoping gives
\begin{equation}
 \label{asm:reconstruction}
 Q_M\alpha_\infty
 =Q_{K-1}\alpha_\infty
  +\sum_{N=K}^MP_N\tau_N
  +\sum_{N=K}^MP_N(\alpha_N-\tau_N)-Z_M.
\end{equation}
The fixed low-frequency term $Q_{K-1}\alpha_\infty$ is a smooth
$L^1$ function by convolution with a Schwartz kernel.  By
Equations~\eqref{asm:two-convergent-series} and
\eqref{asm:variation-tail}, the right side converges in
$L^1+L^2$, hence in distributions, to
\[
 f=Q_{K-1}\alpha_\infty+G+H-Z\in L^1_{\mathrm{loc}}(\R).
\]
The kernels defining $Q_M$ form an approximate identity, so the
left side of Equation~\eqref{asm:reconstruction} converges in
distributions to $\alpha_\infty$.  Therefore
$\alpha_\infty=f\,\cL^1$.  Since $\alpha_\infty$ is a finite
positive measure, this locally integrable density is nonnegative
almost everywhere and belongs to $L^1(\R)$.  This proves absolute
continuity.
\end{proof}

\subsection{Completion of the proof}

\begin{proposition}\label{asm:completion}
For the starting pair in Proposition~\ref{prep:starting-pairs}, there
is a nonzero finite positive measure, absolutely continuous with
respect to $\cL^1$, carried by the original distance set $\Delta(E)$.
Consequently $\cL^1(\Delta(E))>0$.
\end{proposition}

\begin{proof}
Choose the parameters as in Proposition~\ref{asm:shell-decay}.
Lemma~\ref{asm:retained-measures} gives decreasing $\alpha_N$ with
nonzero limit $\alpha_\infty$, and approximants $\tau_N$ with
summable total-variation error.  Proposition~\ref{asm:shell-decay}
gives the second summability hypothesis of
Lemma~\ref{asm:monotone-shell}.  That Lemma makes $\alpha_\infty$
absolutely continuous.  By Lemma~\ref{asm:retained-measures} it is
carried by $\Delta(E)$ and has positive mass, so $\Delta(E)$ cannot
be a Lebesgue null set.
\end{proof}

\begin{proof}[Proof of Theorem~\ref{thm:main}]
Start with a Frostman probability on the given compact set
$E\subset\R^d$, $d\ge2$, with exponent strictly larger than $d/2$.
In the odd-dimensional exceptional affine-plane case,
Lemma~\ref{prep:classical-plane} already proves the conclusion using
original distances inside that plane.  In every other case,
Proposition~\ref{prep:starting-pairs} supplies the separated original
probabilities and filters used above, and
Proposition~\ref{asm:completion} proves the conclusion.  This includes
$d=2$: its cap preparation is the integer case of
Lemma~\ref{prep:integer-caps}.  The auxiliary projections used in
preparation control directions only; the measures in
Equation~\eqref{asm:measures} and their nonzero limit are pushforwards
by $|x-y|$ in $\R^d$.  Thus the argument establishes the asserted
positive measure for the original distance set in every dimension
$d\ge2$.
\end{proof}
